\section{Cas d'études}
\label{sec:reevaluation-base-regles-cas-etude}


Nous présentons deux cas d'études, sous la forme de
constructions, que nous utiliserons tout au long de chapitre afin faciliter la
compréhension des différents mécanismes que nous proposons pour la nomination
persistante des entités topologiques et leurs appariements lors de la
réévaluation. %

\subsection{Construction nº~1}
\label{sec:reevaluation-base-regles-premier-cas}

\todo{Trouver un arrangement pour avoir les évaluations à côté de leurs
	réévaluations}

\begin{figure}[ht!]
	%\renewcommand\thesubfigure{\arabic{subfigure}}
	\centering
	\begin{subfigure}[t]{0.11\textwidth}
		\includegraphics[width=\textwidth]{1-evaluation-square}
		\caption{}
		\label{fig:1-eval}
	\end{subfigure}
	\hfill
	\begin{subfigure}[t]{0.11\textwidth}
		\includegraphics[width=\textwidth]{2-evaluation-extrusion}
		\caption{}
		\label{fig:2-eval}
	\end{subfigure}
	\hfill
	\begin{subfigure}[t]{0.11\textwidth}
		\includegraphics[width=\textwidth]{3-evaluation-insert}
		\caption{}
		\label{fig:3-eval}
	\end{subfigure}
	\hfill
	\begin{subfigure}[t]{0.11\textwidth}
		\includegraphics[width=\textwidth]{4-evaluation-extrusion}
		\caption{}
		\label{fig:4-eval}
	\end{subfigure}
	\hfill
	\begin{subfigure}[t]{0.11\textwidth}
		\includegraphics[width=\textwidth]{5-evaluation-triangulation}
		\caption{}
		\label{fig:5-eval}
	\end{subfigure}
	\hfill
	\begin{subfigure}[t]{0.11\textwidth}
		\includegraphics[width=\textwidth]{6-evaluation-collapse}
		\caption{}
		\label{fig:6-eval}
	\end{subfigure}
	\hfill
	\begin{subfigure}[t]{0.11\textwidth}
		\includegraphics[width=\textwidth]{7-evaluation-chamfer}
		\caption{}
		\label{fig:7-eval}
	\end{subfigure}
	\hfill
	\begin{subfigure}[t]{0.11\textwidth}
		\includegraphics[width=\textwidth]{8-evaluation-colour}
		\caption{}
		\label{fig:8-eval}
	\end{subfigure}
	\caption{Evaluation %'s applications.
		(a) $1$-carré(pos); (b)
		$2$-extrusion($PN_{6}$,vec); (c) $3$-insertion($PN_{46}$); (d)
		$4$-extrusion($PN_{47}$,vec); (e) $5$-triangulation($PN_{33}$); (f)
		$6$-contraction($PN_{63}$); (g) $7$-chanfrein($PN_{55}$); (h) $8$-colorier($PN_{35}$)}
	\label{fig:evaluation}
\end{figure}
%
\begin{figure}[ht!]
	%\renewcommand\thesubfigure{\arabic{subfigure}}
	\centering
	\begin{subfigure}[t]{0.11\textwidth}
		\includegraphics[width=\textwidth]{1-reevaluation-square}
		\caption{}
		\label{fig:1-reeval}
	\end{subfigure}
	\hfill
	\begin{subfigure}[t]{0.11\textwidth}
		\includegraphics[width=\textwidth]{1ADD-reevaluation-insert}
		\caption{}
		\label{fig:1add-reeval}
	\end{subfigure}
	\hfill
	\begin{subfigure}[t]{0.11\textwidth}
		\includegraphics[width=\textwidth]{2-reevaluation-extrusion}
		\caption{}
		\label{fig:2-reeval}
	\end{subfigure}
	\begin{subfigure}[t]{0\textwidth}
		\includegraphics[width=\textwidth]{3-reevaluation-insert}
		\caption{}
		\label{fig:3delete-reeval}
	\end{subfigure}
	\hfill
	\begin{subfigure}[t]{0.11\textwidth}
		\includegraphics[width=\textwidth]{4-reevaluation-extrusion}
		\caption{}
		\label{fig:4-reeval}
	\end{subfigure}
	\hfill
	\begin{subfigure}[t]{0.11\textwidth}
		\includegraphics[width=\textwidth]{5-reevaluation-triangulation}
		\caption{}
		\label{fig:5-reeval}
	\end{subfigure}
	\hfill
	\begin{subfigure}[t]{0.11\textwidth}
		\includegraphics[width=\textwidth]{6-reevaluation-collapse}
		\caption{}
		\label{fig:6-reeval}
	\end{subfigure}
	\hfill
	\begin{subfigure}[t]{0.11\textwidth}
		\includegraphics[width=\textwidth]{7-reevaluation-chamfer}
		\caption{}
		\label{fig:7-reeval}
	\end{subfigure}
	\hfill
	\begin{subfigure}[t]{0.11\textwidth}
		\includegraphics[width=\textwidth]{8-reevaluation-colour}
		\caption{}
		\label{fig:8-reeval}
	\end{subfigure}
	\caption{Reevaluation (a) $1$-carré(pos); (b) ADD$1$-insertion(3); (c)
		$2$-extrusion($PN_{6}$,vec); (d) DELETE $3$-insertion($PN_{46}$); (e)
		$4$-extrusion($PN_{47}$,vec); (f) $5$-triangulation($PN_{33}$); (g)
		$6$-contraction($PN_{63}$); (h) $7$-chanfrein($PN_{55}$); (i) $8$-colorier($PN_{35}$)}
	\label{fig:reevaluation}
\end{figure}

La première construction que nous présentons, illustrée
figure~\ref{fig:evaluation}, est l'objet représenté par la spécification
paramétrique
(liste~\ref{lst:reevaluation-regles-specification-parametrique-spec-init})
ci-après. %

{\footnotesize\ttfamily%
\begin{lstlisting}[%
label=lst:reevaluation-regles-specification-parametrique-spec-init,caption=Spécification
 paramétrique initiale]
1-carre(pos)
2-extrusionFace(PN(6)=[1n6], vec)
3-insertion(PN(46)=[1n5;2n5])
4-extrusionFaceVolume(PN(47)=[1n5;2n6;3n0], vec)
5-triangulationFace(PN(33)=[1n3;2n2])
6-contractionFace(PN(63)=[1n4;2n6;4n6])
7-chanfrein(PN(55)=[1n5;2n6;3n0;4n4;6n2])
8-coloration(PN(35)=[1n3;2n4;5n0])
\end{lstlisting}}

Détaillons cette spécification paramétrique. %
L'application de la première opération, \texttt{carre}, crée une face carrée
(figure~\ref{fig:1-eval}). %
La deuxième opération, \texttt{extrusionFace}, est appliquée sur cette face,
incidente au brin \(6\), et l'extrude en un cube (figure~\ref{fig:2-eval}). %
La troisième opération, \texttt{insertion}, est appliquée sur l'arête
incidente au brin \(46\) et y insère un sommet (figure~\ref{fig:3-eval}). %
La quatrième opération, \texttt{extrusionFaceVolume}, est appliquée sur la
face du dessus, incidente au brin \(47\), et l'extrude en un autre cube
(figure~\ref{fig:4-eval}). %
La cinquième opération, \texttt{trianguleFace}, est appliquée sur la face
avant, incidente au brin \(33\), et la triangule (figure~\ref{fig:5-eval}). %
La sixième opération, \texttt{contractionFace}, est appliquée sur la face
incidente au brin \(63\), et la contracte en un sommet
(figure~\ref{fig:6-eval}). %
La septième opération, \texttt{chanfrein}, est appliquée sur ce même sommet,
incident au brin \(55\), et fait un chanfrein dessus
(figure~\ref{fig:7-eval}). %
Enfin, la dernière opération, \texttt{coloration} colore la face incidente au
brin \(35\) (figure~\ref{fig:8-eval}). %

Nous procédons ensuite à l'édition de cette spécification paramétrique dont la
représentation géométrique, à chaque étape, est illustrée dans la
figure~\ref{fig:reevaluation}. %
Plus précisément, nous ajoutons une opération d'insertion de sommet
(figure~\ref{fig:1add-reeval}) et nous supprimons l'opération
\texttt{3-insertion} afin que celle-ci ne soit pas réévaluée. %

{\footnotesize\ttfamily
\begin{lstlisting}[%
label=lst:reevaluation-regles-specification-parametrique-spec-modif,caption=Spécification paramétrique modifiée]
1-carre(pos)
ADD 9-insertion(3)
2-extrusionFace(PN(6)=[1n6], vec)
DELETE 3-insertion(PN(46)=[1n5;2n4])
4-extrusionFaceVolume(PN(47)=[1n5;2n6;3n0], vec)
5-triangulationFace(PN(33)=[1n3;2n2])
6-contractionFace(PN(63)=[1n4;2n6;4n6])
7-chanfrein(PN(55)=[1n5;2n6;3n0;4n4;6n2])
8-coloration(PN(35)=[1n3;2n4;5n0])
\end{lstlisting}}

  \todo{Décrire les éléments importants de ce cas d'étude et ce qui nous
	intéresse~: scission supplémentaire de la face avant~; fusion par absence de
	création sur le côté~; coloration de deux faces (symétriques) au lieu d'une
	seule}

\subsection{Construction nº~2}
\label{sec:reevaluation-base-regles-deuxieme-cas}

\todo{Générer la deuxième construction} \todo{L'intérêt de cette construction,
  c'est la gestion d'évènements collatéraux pour compléter l'autre cas sans
  l'alourdir~: suppression d'une face à colorier~; fusion de deux faces (dont
  l'une est coloriée) par absence de création}

%%% Local Variables:
%%% mode: LaTeX
%%% TeX-master: "../../../main"
%%% End:
